%----------------------------------------------------------------------------------------
%   USEFUL COMMANDS
%----------------------------------------------------------------------------------------

\newcommand{\dipartimento}{Dipartimento di Matematica ``Tullio Levi-Civita''}

%----------------------------------------------------------------------------------------
% 	USER DATA
%----------------------------------------------------------------------------------------

% Data di approvazione del piano da parte del tutor interno; nel formato GG Mese AAAA
% compilare inserendo al posto di GG 2 cifre per il giorno, e al posto di 
% AAAA 4 cifre per l'anno
\newcommand{\dataApprovazione}{Data}

% Dati dello Studente
\newcommand{\nomeStudente}{Davide}
\newcommand{\cognomeStudente}{Martinelli}
\newcommand{\matricolaStudente}{2077679}
\newcommand{\emailStudente}{davide.martinelli.1@studenti.unipd.it}
\newcommand{\telStudente}{+ 39 373 73 90 678}

% Dati del Tutor Aziendale
\newcommand{\nomeTutorAziendale}{Stefano}
\newcommand{\cognomeTutorAziendale}{Bertolin}
\newcommand{\emailTutorAziendale}{stefano.bertolin@kireygroup.com}
\newcommand{\telTutorAziendale}{}
\newcommand{\ruoloTutorAziendale}{}

% Dati dell'Azienda
\newcommand{\ragioneSocAzienda}{Kirey Group S.r.l.}
\newcommand{\indirizzoAzienda}{Corso Stati Uniti, 14/B, 35127 Padova PD}
\newcommand{\sitoAzienda}{https://www.kireygroup.com/}
\newcommand{\emailAzienda}{}
\newcommand{\partitaIVAAzienda}{}

% Dati del Tutor Interno (Docente)
\newcommand{\titoloTutorInterno}{Prof.}
\newcommand{\nomeTutorInterno}{Tullio}
\newcommand{\cognomeTutorInterno}{Vardanega}

\newcommand{\prospettoSettimanale}{
     % Personalizzare indicando in lista, i vari task settimana per settimana
     % sostituire a XX il totale ore della settimana
    \begin{itemize}
        \item \textbf{Prima Settimana (40 ore)}
        \begin{itemize}
            \item Introduzione al progetto e studio delle tecnologie;
        \end{itemize}
        \item \textbf{Seconda Settimana (40 ore)} 
        \begin{itemize}
            \item Analisi dei requisiti e progettazione tecnica della soluzione da implementare;
        \end{itemize}
        \item \textbf{Terza Settimana (40 ore)} 
        \begin{itemize}
            \item Realizzazione dell'applicativo;
            \item Scrittura documentazione accessoria;
        \end{itemize}
        \item \textbf{Quarta Settimana (40 ore)} 
        \begin{itemize}
            \item Realizzazione dell'applicativo;
            \item Scrittura documentazione accessoria;
        \end{itemize}
        \item \textbf{Quinta Settimana (40 ore)} 
        \begin{itemize}
            \item Realizzazione dell'applicativo;
            \item Scrittura documentazione accessoria;
        \end{itemize}
        \item \textbf{Sesta Settimana (40 ore)} 
        \begin{itemize}
            \item Realizzazione dell'applicativo;
            \item Realizzazione di test dell'applicativo;
            \item Scrittura documentazione accessoria;
        \end{itemize}
        \item \textbf{Settima Settimana (40 ore)} 
        \begin{itemize}
            \item Realizzazione di test dell'applicativo;
            \item Scrittura documentazione accessoria;
            \item Scrittura della documentazione tecnica;
        \end{itemize}
        \item \textbf{Ottava Settimana - Conclusione (40 ore)} 
        \begin{itemize}
            \item Collaudo dell'applicativo;
            \item Conclusione scrittura della documentazione accessoria e tecnica;
        \end{itemize}
    \end{itemize}
}

% Indicare il totale complessivo (deve essere compreso tra le 300 e le 320 ore)
\newcommand{\totaleOre}{320}

\newcommand{\obiettiviObbligatori}{
    \item \underline{\textit{O01}}: Raccolta in tempo reale dei dati di traffico (Logging degli accessi utente, tempi di risposta, ed eventuali errori);

    \item \underline{\textit{O02}}: Individuazione automatica di anomalie con algoritmi rivolti al Machine Learing o AI per identificare:
    \begin{itemize}
        \item Accessi sospetti (es. traffico bot, attacchi DDoS);
        \item Rallentamenti improvvisi del sistema;
        \item Pattern anomali nei comportamenti utente.
    \end{itemize}
    
    \item \underline{\textit{O03}}: Segnalazione automatica delle anomalie rilevate;
    \item \underline{\textit{O04}}: Scrittura di documentazione tecnica e funzionale.
}

\newcommand{\obiettiviDesiderabili}{
	 \item \underline{\textit{D01}}: primo obiettivo;
}

\newcommand{\obiettiviFacoltativi}{
	 \item \underline{\textit{F01}}: primo obiettivo;
}